\documentclass[10pt,a4paper]{amsart}
\usepackage{amsmath,amssymb,amsthm,graphicx}
% For commutative diagrams we use Xy-pic
\usepackage[all]{xy}

% We use 2cell for 2-commutative diagrams.
\xyoption{2cell}
% Theorem environments.
%
\theoremstyle{plain}
\newtheorem{theorem}[subsection]{Theorem}
\newtheorem{proposition}[subsection]{Proposition}
\newtheorem{lemma}[subsection]{Lemma}

\theoremstyle{definition}
\newtheorem{definition}[subsection]{Definition}
\newtheorem{example}[subsection]{Example}
\newtheorem{exercise}[subsection]{Exercise}
\newtheorem{situation}[subsection]{Situation}

\theoremstyle{remark}
\newtheorem{remark}[subsection]{Remark}
\newtheorem{remarks}[subsection]{Remarks}

\numberwithin{equation}{subsection}

% Macros
%
\def\lim{\mathop{\mathrm{lim}}\nolimits}
\def\colim{\mathop{\mathrm{colim}}\nolimits}
\def\Spec{\mathop{\mathrm{Spec}}}
\def\Hom{\mathop{\mathrm{Hom}}\nolimits}
\def\Ext{\mathop{\mathrm{Ext}}\nolimits}
\def\SheafHom{\mathop{\mathcal{H}\!\mathit{om}}\nolimits}
\def\SheafExt{\mathop{\mathcal{E}\!\mathit{xt}}\nolimits}
\def\Sch{\mathit{Sch}}
\def\Mor{\mathop{\mathrm{Mor}}\nolimits}
\def\Ob{\mathop{\mathrm{Ob}}\nolimits}
\def\Sh{\mathop{\mathit{Sh}}\nolimits}
\def\NL{\mathop{N\!L}\nolimits}
\def\CH{\mathop{\mathrm{CH}}\nolimits}
\def\proetale{{pro\text{-}\acute{e}tale}}
\def\etale{{\acute{e}tale}}
\def\QCoh{\mathit{QCoh}}
\def\Ker{\mathop{\mathrm{Ker}}}
\def\Im{\mathop{\mathrm{Im}}}
\def\Coker{\mathop{\mathrm{Coker}}}
\def\Coim{\mathop{\mathrm{Coim}}}

% Boxtimes
%
\DeclareMathSymbol{\boxtimes}{\mathbin}{AMSa}{"02}

%
% Macros for moduli stacks/spaces
%
\def\QCohstack{\mathcal{QC}\!\mathit{oh}}
\def\Cohstack{\mathcal{C}\!\mathit{oh}}
\def\Spacesstack{\mathcal{S}\!\mathit{paces}}
\def\Quotfunctor{\mathrm{Quot}}
\def\Hilbfunctor{\mathrm{Hilb}}
\def\Curvesstack{\mathcal{C}\!\mathit{urves}}
\def\Polarizedstack{\mathcal{P}\!\mathit{olarized}}
\def\Complexesstack{\mathcal{C}\!\mathit{omplexes}}
% \Pic is the operator that assigns to X its picard group, usage \Pic(X)
% \Picardstack_{X/B} denotes the Picard stack of X over B
% \Picardfunctor_{X/B} denotes the Picard functor of X over B
\def\Pic{\mathop{\mathrm{Pic}}\nolimits}
\def\Picardstack{\mathcal{P}\!\mathit{ic}}
\def\Picardfunctor{\mathrm{Pic}}
\def\Deformationcategory{\mathcal{D}\!\mathit{ef}}
\usepackage{cleveref}
\crefname{originalreference}{original reference}{original references}
\Crefname{originalreference}{Original reference}{Original references}
\makeatletter
\newlabel{situation-setup-base-change}{{0FVG}{source}{}{sourcepointer-1}{}}
\newlabel{situation-setup-base-change@cref}{{[originalreference][0][]0FVG}{[1][0][]source}{}{}{}}
\makeatother
\crefname{originallemma}{original lemma}{original lemmas}
\Crefname{originallemma}{Original lemma}{Original lemmas}
\makeatletter
\newlabel{lemma-cycle-push-sheaf}{{02R6}{source}{}{sourcepointer-2}{}}
\newlabel{lemma-cycle-push-sheaf@cref}{{[originallemma][0][]02R6}{[1][0][]source}{}{}{}}
\makeatother
\crefname{originallemma}{original lemma}{original lemmas}
\Crefname{originallemma}{Original lemma}{Original lemmas}
\makeatletter
\newlabel{lemma-pullback-coherent-base-change}{{0FVI}{source}{}{sourcepointer-3}{}}
\newlabel{lemma-pullback-coherent-base-change@cref}{{[originallemma][0][]0FVI}{[1][0][]source}{}{}{}}
\makeatother
\crefname{originallemma}{original lemma}{original lemmas}
\Crefname{originallemma}{Original lemma}{Original lemmas}
\makeatletter
\newlabel{coherent-lemma-flat-base-change-cohomology}{{02KH}{source}{}{sourcepointer-4}{}}
\newlabel{coherent-lemma-flat-base-change-cohomology@cref}{{[originallemma][0][]02KH}{[1][0][]source}{}{}{}}
\makeatother
\title{Proper pushforward of cycle and Chow groups commutes with the base change pullback in the shifted degrees}
\author{The Stacks Project Authors}
\date{stacks-project:chow.tex:local-snapshot-20261012; GNU FDL 1.2 or later}
\begin{document}
\maketitle
\noindent\textit{Editorial note.} All delivered mathematical text is the author's own, copied from the pinned Stacks Project chapter file chow.tex; no statement or proof step was rewritten, reordered or removed.\par\medskip
\noindent\textit{Editorial note.} Preamble provenance: the pinned chapter file opens with an input of the shared preamble and carries no preamble text of its own, so the project's single Stacks-wide preamble staged at sources/stacks-chow/preamble.tex is excerpted verbatim as two spans, lines 14-18 (Xy-pic setup) and lines 154-217 (the native theorem declarations and mathematical macros). Nothing was added, removed or redefined; the wrapper class and the amsmath/amssymb/amsthm/graphicx base packages are supplied by the assembler.\par\medskip
\noindent\textit{Editorial note.} The counted statement and proof contain no cross reference, no citation and no included file, so the delivered body is self contained.\par\medskip
\noindent\textit{Editorial note.} Canonical location of the counted proposition: Stacks tag 0FVL (sources/stacks-chow/pinned-original-source.tex line 15340).\par\medskip
\noindent\textit{Allowed standard prerequisites (editorial).}\begin{itemize}
\item The base change situation with the constant shift c and the function delta'. Same chapter, Situation labelled situation-setup-base-change (source line 15123), invoked in the statement. Fixes the shift c and the dimension conventions in which the two diagrams are stated.
\item The cycle attached to the pushforward of a closed subscheme's structure sheaf. Same chapter, Lemma labelled lemma-cycle-push-sheaf (source line 1992), invoked at proof lines 15366-15368. Supplies the identifications of f\_*[Z] and (f')\_*g\textasciicircum{}*[Z] with the fundamental cycles of the pushed forward structure sheaves.
\item Base change of the k-cycle of a coherent sheaf. Same chapter, Lemma labelled lemma-pullback-coherent-base-change (source line 15193), invoked at proof lines 15369-15370. Supplies g\textasciicircum{}*f\_*[Z] = [(X' -> X)\textasciicircum{}* f\_*O\_Z]\_\{k+c\}; it is counted separately in this lane as S12-algebra-pullback-coherent-base-change.
\item Flat base change for the cohomology of a quasi-coherent sheaf. Cohomology of Schemes, Lemma labelled coherent-lemma-flat-base-change-cohomology (Stacks tag 02KH), invoked at proof lines 15375-15378. Supplies the module isomorphism that closes the argument.
\end{itemize}\medskip
\ifdefined\hypertarget\hypertarget{sourcepointer-1}{}\fi
\noindent\textit{Source reference.} Original reference 0FVG: sources/stacks-chow/pinned-original-source.tex, source lines 15123--15123. Same-chapter pointer into the author environment situation at source line 15123; invoked at proof line None of the counted proof. It is retained only as a navigation pointer; every load-bearing step is inside the delivered span.\par\medskip
\ifdefined\hypertarget\hypertarget{sourcepointer-2}{}\fi
\noindent\textit{Source reference.} Original lemma 02R6: sources/stacks-chow/pinned-original-source.tex, source lines 1992--1992. Same-chapter pointer into the author environment lemma at source line 1992; invoked at proof line 15366 of the counted proof. It is retained only as a navigation pointer; every load-bearing step is inside the delivered span.\par\medskip
\ifdefined\hypertarget\hypertarget{sourcepointer-3}{}\fi
\noindent\textit{Source reference.} Original lemma 0FVI: sources/stacks-chow/pinned-original-source.tex, source lines 15193--15193. Same-chapter pointer into the author environment lemma at source line 15193; invoked at proof line 15369 of the counted proof. It is retained only as a navigation pointer; every load-bearing step is inside the delivered span.\par\medskip
\ifdefined\hypertarget\hypertarget{sourcepointer-4}{}\fi
\noindent\textit{Source reference.} Original lemma 02KH: sources/stacks-chow/tags.txt, source lines 3372--3372. Cross-chapter standard prerequisite located in the project tag list staged in the local snapshot; invoked at proof line 15378 of the counted proof. It is retained only as a navigation pointer; every load-bearing step is inside the delivered span.\par\medskip
\begin{lemma}
\label{lemma-pullback-base-change-pushforward}
In Situation \ref{situation-setup-base-change} let $Y \to X \to S$ be locally
of finite type and let $Y' \to X' \to S'$ be the base change by $S' \to S$.
Assume $f : Y \to X$ is proper. Then $f' : Y' \to X'$ is proper and the diagram
$$
\vcenter{
\xymatrix{
Z_k(Y) \ar[r]_{g^*} \ar[d]_{f_*} & Z_{k + c}(Y') \ar[d]^{f'_*} \\
Z_k(X) \ar[r]^{g^*} & Z_{k + c}(X')
}
}
\quad\text{and}\quad
\vcenter{
\xymatrix{
\CH_k(Y) \ar[r]_{g^*} \ar[d]_{f_*} & \CH_{k + c}(Y') \ar[d]^{f'_*} \\
\CH_k(X) \ar[r]^{g^*} & \CH_{k + c}(X')
}
}
$$
of cycle and chow groups commutes.
\end{lemma}

\begin{proof}
It suffices to show the first diagram commutes. To see this, let
$Z \subset Y$ be an integral closed subscheme of $\delta$-dimension $k$
and denote $Z' \subset X'$ its base change. By construction
we have $g^*[Z] = [Z']_{k + c}$. By Lemma \ref{lemma-cycle-push-sheaf}
we have $(f')_*g^*[Z] = [f'_*\mathcal{O}_{Z'}]_{k + c}$.
By the same lemma we have $f_*[Z] = [f_*\mathcal{O}_Z]_k$. By
Lemma \ref{lemma-pullback-coherent-base-change}
we have $g^*f_*[Z] = [(X' \to X)^*f_*\mathcal{O}_Z]_{k + r}$.
Thus it suffices to show that
$$
(X' \to X)^*f_*\mathcal{O}_Z \cong f'_*\mathcal{O}_{Z'}
$$
as coherent modules on $X'$. As $X' \to X$ is flat and as
$\mathcal{O}_{Z'} = (Y' \to Y)^*\mathcal{O}_Z$, this
follows from flat base change, see
Cohomology of Schemes, Lemma \ref{coherent-lemma-flat-base-change-cohomology}.
\end{proof}

\end{document}
